\section{Current Results and Analysis}

The current evidence answers a narrow part of RQ1: whether the sandbox can
freeze a workload-transfer prediction and compare it with stable, previously
unseen measurements without leakage.  It does not yet answer RQ2--RQ5.

\subsection{R1a: End-to-End Blind Workload Transfer}

All three calibration and three holdout runs complete without failed requests.
Across the four primary holdout metrics, coefficients of variation range from
\PilotCVRange{}, substantially below the prediction error.  The frozen L0
prediction and independent L1 measurements appear in
Table~\ref{tab:pilot-validation}.

\begin{table}[t]
\centering
\caption{Protocol-frozen blind workload-transfer pilot on one H100.  The L0
prediction was produced before measurement; the measured column is the median
$\pm$ MAD over three L1 repeats.  PI indicates whether the pre-frozen prediction
interval contains the observed median.}
\label{tab:pilot-validation}
\scriptsize
\setlength{\tabcolsep}{2.5pt}
\begin{tabular}{@{}lrrrr@{}}
\toprule
Metric & L0 pred. & L1 measured & APE & PI \\
\midrule
Energy (J/1k out.) & \PilotPredEnergy & $\PilotObsEnergy\!\pm\!\PilotMADEnergy$ & \PilotEnergyAPE & \cmark \\
Throughput (tok/s) & \PilotPredThroughput & $\PilotObsThroughput\!\pm\!\PilotMADThroughput$ & \PilotThroughputAPE & \cmark \\
P95 TTFT (ms) & \PilotPredTTFT & $\PilotObsTTFT\!\pm\!\PilotMADTTFT$ & \PilotTTFTAPE & \cmark \\
P95 TPOT (ms) & \PilotPredTPOT & $\PilotObsTPOT\!\pm\!\PilotMADTPOT$ & \PilotTPOTAPE & \cmark \\
\midrule
Four-metric summary & -- & -- & \PilotMeanMAPE & \PilotCoverage \\
\bottomrule
\end{tabular}
\end{table}


Point-prediction MAPE is \PilotMeanMAPE{}.  Energy and TPOT are overpredicted by
\PilotEnergyAPE{} and \PilotTPOTAPE{}, while throughput and TTFT are
underpredicted by \PilotThroughputAPE{} and \PilotTTFTAPE{}.  The largest error
is therefore the paper's primary efficiency metric, not a secondary diagnostic:
one anchor is insufficient to rank nearby energy candidates.

\subsection{Protocol Integrity and Interval Interpretation}

The manifest confirms that prediction precedes measurement and that all repeats
are comparable.  All \PilotCoverage{} medians lie inside their intervals, but
the profile assigns \PilotRelativeUncertainty{} relative uncertainty and has
not seen an independent validation set.  This is diagnostic containment, not
calibrated coverage: one workload and broad intervals cannot establish nominal
coverage or support a Pareto decision.

\subsection{Next Frozen Validation and Remaining Claims}

The frozen six-point campaign next estimates workload-transfer error and
recalibrates interval width before an untouched holdout.  RQ2--RQ5 still require
multi-configuration frontiers, L2--L4 evidence, matched-budget baselines,
cross-hardware transfer, agent ablations, and target-scale reruns.  We therefore
make no current claim about Pareto recall, GPU-hour reduction, agent advantage,
or deployment energy savings.
